\chapter{The Thorin phase diagram}\label{chap:phase}

This chapter sets out the most promising computational path to RH that the book has found, and says exactly where it stops. The chapters before this one compute the Thorin measure of the clock in many ways and at many resolutions, and Section~\ref{ch:finite} explains why none of those computations, by itself, can settle RH: a perturbation of $\xi$ with the same zeros below a large height passes every test that sees only that height. This chapter asks the question that the result leaves open. Which computations see past the verified height at all, and of those, which comes closest to RH?

The answer has three parts. First, only three routes in the book turn a finite computation into a statement about all zeros, and each does so because some analytic argument takes over above a barrier: the Euler product at basepoints $\alpha\ge1$ (Section~\ref{ch:basepoint}), the functional equation for the symmetrised truncations of Section~\ref{ch:rz4}, and the heat flow at positive time (Section~\ref{ch:dbn}). Second, the third of these is the Thorin programme itself, carried out with a second parameter. The shifted clock of the time $t$ function $\xi_t$ at basepoint $\alpha$ is a generalised gamma convolution exactly when no zero of $\xi_t$ lies to the right of $\alpha$, and in this language de Bruijn's strip theorem says that the set of points $(t,\alpha)$ where it is one is carried along parabolas by the flow and that $\Lambda$ is the minimum of $t+2(\alpha-\frac12)^2$ over it. RH is the corner $(0,\frac12)$. Third, the Thorin density of $\xi_{0.2}$, computed from the P\'olya integral with no primes and no zeros, is close to the function that the same computation gives for $\xi$ (a Thorin density only if the clock of $\xi$ at that basepoint is a GGC) on the windows computed in Section~\ref{cp:sec:see}, below height $100$: at the basepoint $0.6$ its valleys change by a median of $2.3\%$ and at most $14.2\%$. What separates the theorems $\Lambda\le0.22$ of Polymath15 \cite{Polymath} and $\Lambda\le0.2$ of Platt and Trudgian \cite{PT21} from RH is not anything a computed density shows; it is the existence of the analytic tail at $t=0.2$ and its absence at $t=0$.

So the chapter answers a question that the reader may have asked since Chapter~\ref{chap:measure}: whether the sine squared lemma and the Thorin side offer a computational alternative to the de Bruijn Newman flow. They do not offer an alternative; they offer the right language for it. The Thorin side is the object to compute, prime free and well defined at every basepoint, and the flow is the only known source of an analytic tail below the basepoint $\frac78$.

Section~\ref{cp:sec:routes} classifies the routes of the book by what lies above their barrier. Section~\ref{cp:sec:transfer} proves the threshold theorem for the clocks of $\xi_t$ at every basepoint. Section~\ref{cp:sec:phase} states and proves the phase diagram theorem and places the known results on it. Section~\ref{cp:sec:see} computes the Thorin density of $\xi_t$ at two times and two basepoints. Section~\ref{cp:sec:cert} writes the certificate for one point of the diagram with explicit constants and carries out its finite part in ball arithmetic below height $160$. Section~\ref{cp:sec:uniform} reports two experiments on where the missing uniform step could come from, a Weil Gram matrix on compactly supported test functions and the zero gas along the flow, fits the cost of the flow route to the bound of Polymath15 and their table of conditional bounds, and sets out a programme and open problems.

Sections~\ref{cp:sec:transfer}, \ref{cp:sec:cert} and the two experiments of Section~\ref{cp:sec:uniform} were first written as a working report by the authors \cite{FA26}, which this chapter supersedes. Results are marked as in the rest of the book: a result with a proof in the text is proved, with no hypothesis unless one is stated, and a number in a table is computed by the script named in the caption, one of the programs \texttt{phase\_*.py} (Appendix~\ref{app:code}). We use the normalisation of Section~\ref{ch:dbn}: $H_t(z)=\int_0^\infty e^{tu^2}\Phi_{\rm dBN}(u)\cos(zu)\,du$, $\xi_t(s)=8H_t(-2i(s-\frac12))$, a zero $x+iy$ of $H_t$ is the zero $\frac12-\frac y2+\frac{ix}2$ of $\xi_t$, and $P_a(x)=a/(\pi(a^2+x^2))$ is the Cauchy density of scale $a>0$.
